%---------------------------------------------------------------
%  SECTION III — SYSTEM ARCHITECTURE AND PROTOCOL DESIGN
%---------------------------------------------------------------
\section{System Architecture and Protocol Design}
\label{sec:architecture}

\subsection{Four-Tier Architecture}

Q-SFTP is structured as a four-tier modular system, as shown in
Fig.~\ref{fig:system_arch}.

\begin{figure}[htbp]
  \centering
  \includegraphics[width=\columnwidth]{system_architecture}
  \caption{Q-SFTP four-tier system architecture.}
  \label{fig:system_arch}
\end{figure}

\begin{enumerate}
  \item \textbf{Client Layer}: A single-page application (HTML5/CSS3/
        JavaScript) running in the browser. Performs client-side
        BLAKE2b-256 hashing via \texttt{blakejs} and persists transfer
        state in \texttt{localStorage}.
  \item \textbf{Application Layer}: Flask~3.x on port~5000 handles
        REST API endpoints, user sessions, integrity coordination,
        privacy management, and admin panel.
  \item \textbf{Transport Layer}: PQC handshake server on TCP port~8888
        implements the Q-TLS protocol. All file data flows through this
        fully encrypted channel.
  \item \textbf{Storage Layer}: Per-user isolated directories
        (\texttt{ServerStorage/\{username\}/}) and a shared folder
        (\texttt{ServerStorage/shared/}) for collaborative access.
\end{enumerate}

\subsection{Q-TLS Handshake Protocol}

The Q-TLS handshake achieves mutual authentication and establishes a
forward-secret session key in four phases, as illustrated in
Fig.~\ref{fig:pqc_handshake}.

\begin{figure}[htbp]
  \centering
  \includegraphics[width=\columnwidth]{pqc_handshake}
  \caption{Four-phase Q-TLS post-quantum handshake.}
  \label{fig:pqc_handshake}
\end{figure}

\textbf{Phase 1 — Hello Exchange.}
The client sends a \textit{ClientHello} containing its Dilithium2
certificate, a random nonce $N_C$, and a UTC timestamp $T_C$. The
server verifies the certificate against the CA's pinned public key,
checks that $|T_S - T_C| < 30$~s (replay protection), generates a
fresh Kyber512 key pair $(pk_K, sk_K)$, and replies with a signed
\textit{ServerHello}:
\[
  \sigma_S = \text{Dilithium2.Sign}(sk_S,\, N_C \| N_S \| pk_K)
\]

\textbf{Phase 2 — Kyber512 Key Encapsulation.}
The client verifies $\sigma_S$, then encapsulates a shared secret:
\[
  (ct, ss) = \text{Kyber512.Encaps}(pk_K)
\]
The ciphertext $ct$ (768~B) is sent to the server, which recovers
$ss$ via $\text{Kyber512.Decaps}(sk_K, ct)$.

\textbf{Phase 3 — Mutual Authentication.}
Both parties compute a transcript hash:
\[
  h_{tr} = \text{SHA-256}(\text{CH} \| \text{SH} \| \text{CKS})
\]
The client signs $h_{tr}$ with Dilithium2; the server verifies it,
completing \emph{mutual} authentication.

\textbf{Phase 4 — Secure Channel.}
The session key is derived via HKDF-SHA256:
\[
  K = \text{HKDF}(ss,\; \text{salt}{=}\text{SHA-256}(N_C\|N_S),\;
      \text{info}{=}\texttt{"Q-SFTP-SESSION-KEY"})
\]
All subsequent communication uses AES-256-GCM with per-message
96-bit nonces.

\subsection{BLAKE2b-256 File Integrity}

Every file undergoes dual-side hash verification:
\begin{enumerate}
  \item Client computes $H_C = \text{BLAKE2b-256}(F)$ before upload.
  \item Server computes $H_S = \text{BLAKE2b-256}(F)$ after reception.
  \item Mismatch ($H_C \neq H_S$) signals in-transit corruption.
  \item All hashes are persisted in \texttt{hash\_registry.json} and
        periodically re-verified by a background \texttt{IntegrityChecker}
        thread to detect at-rest tampering.
\end{enumerate}

\subsection{Resumable Chunked Transfer Protocol}

Figure~\ref{fig:resumable} shows the chunked transfer design.

\begin{figure}[htbp]
  \centering
  \includegraphics[width=\columnwidth]{resumable_transfers}
  \caption{Chunked resumable transfer with persistence.}
  \label{fig:resumable}
\end{figure}

Files are split into 256~KB chunks. Each chunk is independently
AES-256-GCM encrypted and transmitted. The server maintains a
\texttt{transfer\_log.json} with the current upload offset; the client
mirrors this in \texttt{localStorage}. On reconnection, the server
returns the saved \texttt{resume\_offset}, and the client resumes from
the last confirmed byte—surviving browser refreshes, connection drops,
and server restarts.

\subsection{Role-Based Access Control (RBAC)}

Q-SFTP enforces a three-tier RBAC model, as depicted in
Fig.~\ref{fig:rbac}:

\begin{figure}[htbp]
  \centering
  \includegraphics[width=\columnwidth]{rbac}
  \caption{Three-tier RBAC permission model with directory isolation.}
  \label{fig:rbac}
\end{figure}

\begin{table}[htbp]
  \centering
  \caption{RBAC Permission Matrix}
  \label{tab:rbac}
  \begin{tabular}{lccc}
    \toprule
    \textbf{Operation} & \textbf{Admin} & \textbf{Standard} & \textbf{Guest} \\
    \midrule
    List / Download  & \checkmark & \checkmark & \checkmark \\
    Upload / Mkdir   & \checkmark & \checkmark & $\times$ \\
    Delete           & \checkmark & \checkmark & $\times$ \\
    Admin Panel      & \checkmark & $\times$   & $\times$ \\
    \bottomrule
  \end{tabular}
\end{table}

\subsection{Privacy and File Protection}

\textbf{Metadata stripping:} The \texttt{PrivacyManager} removes
EXIF data from JPEG/PNG images (Pillow), author/date metadata from
PDFs (PyPDF2), and core properties from DOCX files (python-docx)
before storage.

\textbf{File validation:} \texttt{FileValidator} blocks 17 dangerous
extensions and verifies file magic numbers for 20+ binary formats,
preventing extension-spoofing uploadsof malware.

\textbf{Audit logging:} All user actions are stored in SQLite with IP
anonymisation (SHA-256 with a daily rotating salt) and log-chain
hash integrity.
